\cotefile{75}
\scanpage{26}
c) Orientation d'un ensemble fini $I$\inserted{(cardinal $n$).} La donnée de $I$ équivaut à celle d'un torseur
\[
 P=\operatorname{Isom}(\deleted{\unsure{[1,n]_{\mathbb Z}|[1,n]_{\mathrm{ent}}}}I_n,I)=\unsure{\operatorname{numérot}|\operatorname{numérot.}}(I)
\]
$(I_n\mathrel{\unsure{=|:}}\{m\in\mathbb N^{*}\mid m\leq n\}=\{\deleted{\unsure{0}}1,\unsure{2},\ldots,n\})$ sous le groupe
\[
 \deleted{\unsure{\mathfrak S|S}}\mathfrak S_n=\mathfrak S_{I_n}=\operatorname{Aut}(I_n).
\]

On a un hom canonique ``signature''
\diagram{\begin{tikzcd}[ampersand replacement=\&]
\mathfrak S_n \arrow[r,"s=\mathrm{sgn}"] \& \{\pm1\} \&[-1em] \text{noyau }\mathfrak A_n
\end{tikzcd}}
et par définition l'\unsure{ens} des 2 orientations de $I$ est
\[
 \unsure{\omega|w}_I=\deleted{\unsure{\operatorname{Isom}(|\operatorname{Isom}}}P\times_{\mathfrak S_n}\{\pm1\}\simeq P/\mathfrak A_n .
\]

\uline{Compatibilités avec} a) \uline{et} b)

\uline{Avec} b) On a sur $E_{(I,k)}=k^{I}$ un produit scalaire canonique
\[
 (x,y)=\sum_{i\in I}x_i y_i
\]
pour lequel $(e_i)$ est une base orthonormale. Ainsi $\deleted{\unsure{E}}I$ \deleted{\unsure{donne}}\inserted{\unsure{donne}} naissance à un module quadratique $E(I,\unsure{x|k|X})$. Ceci posé, on \inserted{définit} une \unsure{bijection} canonique
\[
 \omega_{E(I,x)}\simeq\omega_I
\]
\margin{Si $i_1,\ldots,i_n$ est une \unsure{numérotation|numération} de $I$, $e_{i_1}\ \unsure{\cdots|\text{--}}\ e_{i_{\unsure{n}}}$ est une base \inserted{\unsure{ord.}} \unsure{orthon.} de $E$, \ldots}
par transitivité dans
\diagram{\begin{tikzcd}[ampersand replacement=\&,column sep=large]
 \mathfrak S_n \arrow[r,hook,"\mathrm{incl.}"] \arrow[rr,bend right=24,"\mathrm{sgn}" near start] \& \mathrm{O}(\unsure{n|N},k) \arrow[r,"\det"] \& \{\pm1\}
\end{tikzcd}}

En particulier, on trouve (\unsure{comp.} avec a).),
\[
 \omega_{E(I,\mathbb R)}\simeq\omega_I
\]
% Arrow from the preceding low E(I,R) toward the note along the foot of the sheet.
\margin{\tikz[baseline=-.35ex]{\draw[->] (0,.12) .. controls (.17,.08) and (.31,-.06) .. (.49,-.08);}\,\unsure{considéré} comme \unsure{variété} \unsure{via|vide} \unsure{ses pts|ses points|supra}}
% archive sheet number: 25

\scanpage{27}
\uline{Formules d'\unsure{associativité|associativités}}

a) Vectoriels réels. \unsure{Soient|Soit} $I$ ens fini, $(E_i)_{i\in I}$ famille de vectoriels \inserted{réels} de dim finie, $E=\coprod_{i\in I}E_i$ la somme directe ($\simeq$ produit direct). \deleted{\unsure{Soit}} $I^{+}$ $(I^{-})$ l'ens des $i\in I$ tels que $d_i=\operatorname{rg}_{\mathbb R}E_i$ soit pair (impair). On \unsure{a un} iso\unsure{m} canonique
\[
 \boxed{\unsure{\omega|w}_E\simeq\omega_{I^{-}}\wedge\bigwedge_{i\in I}\omega_{E_i}}
\]
obtenu ainsi, en regardant le $2^{\text{ième}}$ mb comme \deleted{\unsure{un}} quotient de
{\reversemarginpar\margin{\figure{\begin{tikzpicture}
\node[inner sep=2pt] (nb) {\parbox{.86\linewidth}{\uline{NB} Un ordre total de $I$ en induit un de $I^{-}$, donc
\diagram{\resizebox{\linewidth}{!}{\begin{tikzcd}[ampersand replacement=\&,column sep=small]
\operatorname{numérot}(I)\arrow[r,"\text{épi}"] \& \operatorname{numérot}(I^{-})
\end{tikzcd}}}}};
\draw ([xshift=3pt]nb.north west) .. controls (nb.north) and ([xshift=-3pt]nb.north east) .. (nb.north east)
 .. controls ([xshift=-2pt]nb.east) and (nb.east) .. ([xshift=2pt]nb.south east)
 .. controls (nb.south east) and (nb.south) .. ([xshift=3pt]nb.south west);
\end{tikzpicture}}}\normalmarginpar}
\[
 \operatorname{numérot}(I^{-})\inserted{\times\operatorname{numérot}(I^{+})}\times\prod_{i\in I}\operatorname{bases}(E_i)
\]

Soit $\{i_1,\ldots,i_{n^{-}}\}$ \unsure{num.} de $I^{-}$, \inserted{$j_1\ \cdots\ j_{n^{+}}$ une autre de $I^{+}$,} et $\unsure{\forall}\ i\in I$, $B_{\unsure{i|\alpha}}=(\deleted{\unsure{B}}e_{i,1},\ldots,e_{i,d_i})$ une base ordonnée de $E_i$, alors
\[
 e_{i_1,1}\unsure{2|,},\ldots,e_{i_1,d_1},e_{i_2,\unsure{1|2}},\ldots e_{i_2,d_2},\ldots,e_{i_n,1},\ldots,e_{i_n,d_{i_n}}
\]
est une base de $E$. Si on change une $B_i$ par une matrice de $\mathrm{Gl}(d_i,\mathbb R)^{+}$, alors la base se trouve changée par une matrice de $\mathrm{Gl}(d,\mathbb R)^{+}$ $(d=\sum d_i\deleted{\unsure{)}}=\operatorname{rg}E)$, \inserted{donc définit \unsure{=|même} él. de $\omega_E$.} Si on change l'ordre de $I^{-}$ par une permutation $\sigma^{-}$ (une transposition entre $\alpha$ et $\alpha+1$, disons \ldots) On change la base par une matrice de dét $\deleted{\unsure{(-1)^{i_{\alpha}i_{\alpha+1}}|(-1)^{i_{\alpha}+i_{\alpha+1}}}}\operatorname{sgn}\sigma^{-}$ ; par contre, si on change \deleted{\unsure{par}} l'ordre de $I^{+}$ par une $\sigma^{+}$, la matrice de passage est de det $+1$ en tous cas \ldots
% archive sheet number: 26

\scanpage{28}
b) Contexte des familles \unsure{finies} de modules quadratiques : $=$ formule, avec $=$ \unsure{démonstration|démonstrati.}

c) \diagram{\begin{tikzcd}[ampersand replacement=\&]
\deleted{\unsure{I}}J \arrow[r,"p"] \& I
\end{tikzcd}}
application \inserted{(surjection \unsure{si on veut})} d'ens \deleted{\unsure{C|f}}finis,
\[
 I^{-}=\{i\in I\mid\operatorname{Card}p^{-1}(i)\ \text{impair}\}.
\]
\deleted{\unsure{Je dis que}} [p. ex $J$ ens avec relation d'équivalence \inserted{$\mathcal R$}, et $I=J/\mathcal R$]. \deleted{\unsure{O}}Alors
\[
 \unsure{\omega|w}_J\simeq\omega_{I^{-}}\wedge\bigwedge_{i\in I}\omega_{J_{\unsure{i|2}}}
\]
% slanted
\margin{\illegible{1}}

\uline{Compatibilités} entre les \unsure{isom|isom.|trois} \unsure{can.|cas.} a) b) c)

\figure{\begin{tikzpicture}
\draw (0,0) .. controls (.32,.03) and (.92,-.04) .. (1.2,-.01);
\end{tikzpicture}}

\uline{Cas d'extensions}
\[
 \{\unsure{0|\cdot}\}=\deleted{\unsure{0}}E_0\subset E_1\subset E_2\ \cdots\ \subset E_d=E
\]
\deleted{\unsure{On}} Définir un \unsure{isom.} canonique
\[
 \omega_E\simeq\bigwedge_{1\leq i\leq d}\omega_{E_i\unsure{/|{}}E_{i-1}}
\]
\figure{\begin{tikzpicture}
\draw (0,0) .. controls (.76,.06) and (1.95,.16) .. (2.90,.17);
\draw (1.04,-.05) .. controls (1.50,-.04) and (2.10,-.09) .. (2.46,-.05);
\end{tikzpicture}}
% archive sheet number: 27

\scanpage{30}
\figure{\begin{tikzpicture}
\node[anchor=north west,inner sep=0pt,outer sep=0pt] (letter) at (0,0) {\begin{minipage}{0.98\linewidth}
\footnotesize
% letterhead
\begin{minipage}[t]{0.56\linewidth}
\centering\sffamily
UNIVERSITE DES SCIENCES ET TECHNIQUES\\
DU LANGUEDOC\\
\rule{0.55cm}{0.4pt}\\
INSTITUT DE MATHEMATIQUES\\
Place Eugène Bataillon\\
\unsure{34060|34050} MONTPELLIER CEDEX\\
\rule{0.65cm}{0.4pt}\\
\end{minipage}\hfill
\begin{minipage}[t]{0.40\linewidth}
Montpellier, le\\[2.4\baselineskip]
Villecun le 30.01.78\\
\end{minipage}

\smallskip
Monsieur Alexandre Grothendieck\\
Professeur\\

\hfill\begin{minipage}{0.49\linewidth}
à Monsieur le Ministre de\\
l'Education Nationale\\
\end{minipage}

\smallskip
Monsieurle Ministre,\\

\deleted{Je me permets de vous envoyer copie d'une lettre datée de ce même jour, adressée à M. Rouzaud, Administrateur Provisoire de l'U.S.T.L., dans laquelle je lui soumets certains des problèmes de fonctionnement posés aux Universités par l'article 21 de l'ordonnance du 2.11.1945.}

\smallskip
\deleted{Par le même courrier j'envoie également copie de la même lettre à M. Peyrefitte, Ministre de la Justice, en attirant son attention sur les difficultés créées par cette ordonnance. Celles-ci ne touchent certes pas, à mon sens, l'aspect le plus important des problèmes soulevés par l'existence même ce cette ordonnance. Par contre, par leur nature technique, elles me semblent bien du ressort des responsabilités' officielles d'un Administrateur ou d'un Ministre. C'est pourquoi il m'a semblé opportu\deleted{\unsure{m|n}}\unsure{n|m} de vous \makebox[0pt][l]{\raisebox{0.65ex}{\tiny\unsure{ç|c}}}soumettre, en votre qualité de Ministre de l'Education Nationale, des problèmes de votre compétence. Il me semblerait extrêmement souhaitable que vous preniez position publiquement à leur sujet.}

\smallskip
\hspace*{0.18\linewidth}\deleted{Veuillez agréer, Monsieur le Ministre, l'expression de}\\
\deleted{ma considération très distinguée .}\\

\hfill\handwritten{A Grothendieck}\\
\hspace*{0.69\linewidth}\handwritten{\unsure{''}}\\
\end{minipage}};
\begin{scope}[x={(letter.north east)},y={(letter.south west)}]
\draw[gray,very thin] (0.654,0.258) -- (0.644,0.235)
 .. controls (0.561,0.272) and (0.327,0.389) .. (0.401,0.404)
 .. controls (0.449,0.420) and (0.738,0.382) .. (0.651,0.456)
 .. controls (0.579,0.511) and (0.399,0.568) .. (0.422,0.624)
 .. controls (0.451,0.668) and (0.719,0.641) .. (0.654,0.713)
 .. controls (0.593,0.781) and (0.477,0.842) .. (0.508,0.994);
\draw[gray,very thin] (0.654,0.254)
 .. controls (0.588,0.297) and (0.461,0.379) .. (0.443,0.407)
 .. controls (0.414,0.445) and (0.537,0.422) .. (0.528,0.449)
 .. controls (0.517,0.477) and (0.402,0.533) .. (0.492,0.583)
 .. controls (0.590,0.638) and (0.508,0.674) .. (0.475,0.733)
 .. controls (0.428,0.806) and (0.414,0.817) .. (0.431,0.889);
\end{scope}
\end{tikzpicture}}
% archive sheet number: 29

\scanpage{31}
% letterhead
\noindent\begin{minipage}[t]{0.57\linewidth}
\centering\sffamily
UNIVERSITE DES SCIENCES ET TECHNIQUES\\
DU LANGUEDOC\\
\rule{0.55cm}{0.4pt}\\
INSTITUT DE MATHEMATIQUES\\
Place Eugène Bataillon\\
34060 MONTPELLIER CEDEX\\
\rule{0.65cm}{0.4pt}
\end{minipage}\hfill
\begin{minipage}[t]{0.37\linewidth}
Montpellier, le
\end{minipage}

\begin{center}
\setlength{\ULdepth}{1.45\baselineskip}
\uline{COURS\quad C4 1977-78 de $\mathbb A$. GROTHENDIECK}\par\vspace{.6\baselineskip}
\end{center}

\uline{Géométrie de l'icosaèdre}

Il est étrange que l'icosaèdre, étudié depuis plus de deux mille ans dans des contextes et des optiques très divers, ne semble jamais avoir été étudié systématiquement du point de vue combinatoire, ni de celui des réalisations géométriques sur des corps (voire des anneaux) généraux. Ce ``cours'' se voudrait une investigation commune sur ce thème, en partant de l'intuition combinatoire la plus naïve, exemplifiée sur des modèles en carton (avec un coloriage convenable des arètes).

Nous aborderons donc la théorie de l'i\deleted{\unsure{e}}cosaèdre d'un point de vue combinatoire d'abord, en explicitant dans le détail le dictionnaire entre la géométrie combinatoire de l'icosaèdre orienté, et celle de l'ensemble ``orienté'' à cinq éléments (les cinq couleurs coloriant les arètes !) - ce qui donnera en même temps une motivation et une inspiration géométrique pour étudier le groupe des automorphismes commun aux deux situations, savoir le groupe alterné $\handwritten{\unsure{\mathfrak A|\mathcal A}}_5$ (qui est aussi le plus petit groupe simple non abélien). La question de construire des ``réalisations géométriques'' de l'icosaèdre combinatoire, comme ``polyèdre régulier'' sur le corps des réels, ou sur tout autre corps, sera l'occasion à la fois d'approfondir la compréhension de la structure combinatoire, et de nous familiariser avec les notions de géométrie vectorielle et affine (voire de géométrie projective) sur des corps de caractéristique quelconque.
% archive sheet number: 30

\scanpage{32}
% In each composite field letter the extra stroke is marked on \mathbb. The F and its lowered numeral are typewritten.
En relation avec l'icosaèdre, nous découvrirons que les corps de caractéristique deux et cinq jouent un rôle très particulier, tant pour l'étude des réalisations géométriques, que pour l'étude combinatoire abstraite, qui peut s'expliciter par des notions de géométrie sur les corps ${\handwritten{\mathbb}{F}}_4$ et ${\handwritten{\mathbb}{F}}_5$ - ces deux dictionnaires étant subordonnés aux ``isomorphismes exceptionnels'' :
\[
 \begin{aligned}
 \handwritten{\unsure{\mathfrak A|\mathcal A}}_5 &\simeq \mathrm{Sl}\,(2,{\handwritten{\mathbb}{F}}_4)\quad(\simeq\mathrm{GP}(1,{\handwritten{\mathbb}{F}}_4)\simeq\mathrm{SO}(3,{\handwritten{\mathbb}{F}}_4))\\
 \handwritten{\mathfrak S}_5 &\simeq \mathrm{GP}\,(1,{\handwritten{\mathbb}{F}}_5)\quad(\simeq\mathrm{SO}(3,{\handwritten{\mathbb}{F}}_5))\ .
 \end{aligned}
\]
% archive sheet number: 31

\scanpage{33}
\begin{center}
Cours DEA 1977/78 de A. Grothendieck

\uline{Géométrie arithmétique des polygônes réguliers}
\end{center}

{\reversemarginpar\margin{\handwritten{!}}\normalmarginpar}%
% A short hand-drawn line under the typed words ``en commun''.
Ce ``cours'' se fera dans l'esprit d'une investigation \handwritten{\uline}{en commun} de multiples facettes de la notion de polygône régulier au-dessus d'un corps de base, puis d'un anneau de base quelconque. Ce sera une occasion stimulante pour nous familiariser avec les notions de base de géométrie vectorielle, affine, projective - avec ou sans métrique - en dimensions un, deux et trois, et avec les phénomènes particuliers à ces dimensions, ainsi qu'avec les aspects parfois imprévus que prennent des notions familières (comme celle de réflexion\deleted{\unsure{s}}, ou de centre d'un polygône régulier) sur les corps de caractéristique positive, et notamment sur les corps finis. Le besoin d'une géométrie sur des anneaux de base généraux (tel l'anneau ``sur lequel s'écrit'' le $\mathrm{n}$-polygône régulier ``le plus général'') se fera sentir en cours de route, et nous amènera sans doute à affleurer quelques points de fondements de géométrie algébrique. Si le temps le permet et si d'autres investigations ne nous retiennent , nous tenterons, dans la volée, de construire l'icosaèdre régulier ``général'' sur un anneau de base convenable, fini sur l'anneau ${\handwritten{\mathbb}{Z}}$ des entiers, ayant des corps résiduels de toutes caractéristiques. (La caractéristique deux \hspace{.5em}donnera sans doute du fil à retordre !).

A partir d'un point de départ intuitif et familier, ce ``cours'' voudrait être une ouverture à un univers fascinant, pleins de mystères qui ne demandent qu'à être élucidés. Il ne peut donner des fruits qu'avec une participation active de l'étudiant et un travail personnel important, notamment pour développer a\deleted{\unsure{o|e}}vec le détail nécessaire les nombreux points \handwritten{(}de fondement\inserted{\handwritten{s}} notamment) qui n'auront été qu'esquissés dans le ``cours''. Il lui est très fortement conseillé de suivre en même temps le cours C4 sur l'icosaèdre, qui se fera dans un esprit analogue. Il y aura sans doute de nombreuses interférences entre ces deux cours de réflexion.
% archive sheet number: 32
